\documentclass[11pt,a4paper]{article}

\usepackage[utf8]{inputenc}
\usepackage[T1]{fontenc}
\usepackage[english]{babel}
\usepackage{amsthm}
\usepackage{geometry}
\usepackage{enumitem}
\usepackage{setspace}
\usepackage{parskip}
\usepackage{amsmath}
\usepackage{amssymb}
\usepackage{booktabs}
\usepackage{array}
\usepackage{microtype}
\usepackage{hyperref}

\geometry{margin=2.5cm}
\onehalfspacing

\theoremstyle{definition}
\newtheorem{definition}{Definition}[section]
\newtheorem{claim}{Claim}[section]
\newtheorem{thesis}{Thesis}[section]
\newtheorem{proposition}{Proposition}[section]

\theoremstyle{remark}
\newtheorem{remark}{Remark}[section]
\newtheorem{objection}{Objection}[section]
\newtheorem{example}{Example}[section]

\title{\textbf{Information Idealism}\\[0.3em]
\large An Informational, Axiologically Neutral, and Domain-Absolutely
Independent Ontology of Eidos}
\author{Gasan Bagirov}
\date{September 2026}

\begin{document}
\maketitle

\begin{abstract}
\noindent
This document formalizes a reinterpretation of Plato's theory of Forms
(eidos). On this view, an \emph{Idea} is not a living being, not a moral
entity, and not a construct of human consciousness. It is
\emph{information}: an objective, eternal, immutable, and aseitic
structure. Truth, falsehood, good, and evil are not properties of Ideas
themselves but \emph{human markers} applied to our cognitive relation
with them. The position preserves the realist core of Platonism while
discarding its vitalistic and moralistic commitments; we call it
\textbf{Information Idealism}. The formal core is a category-theoretic
account of structures, their invariance, and the adequacy of
representations, supplemented by an explicit theory of
misrepresentation as failure of injectivity on blocks. Two features of
the position deserve emphasis at the outset. First, the independence
claimed for Ideas is \emph{domain-absolute}: it excludes dependence on
any mind, but not dependence on the constitutive transformations of a
domain (Definition~\ref{def:domain-abs}). Second, the theory of
normative markers is \emph{constructivist} in the sense of
Sellars~\cite{Sellars1956} and Brandom~\cite{Brandom1994}: the
normative force of ``true'' and ``good'' is located in the pragmatic
commitments of agents, not derived from the ontology of structures. The
position is programmatic: several transitions are argued
heuristically, and this is flagged where it occurs. A worked example
(\S\ref{sec:example}) and a section of open problems
(\S\ref{sec:limitations}) are included to make the position
assessable.
\end{abstract}

\section{Introduction}

Classical Platonism holds that the eidos is:
\begin{itemize}[nosep]
  \item eternal and unchanging,
  \item the source of being and knowledge,
  \item subordinate to the Idea of the Good,
  \item accessible to the soul through recollection.
\end{itemize}

Information Idealism preserves the \emph{ontological realism} of
Platonism but rejects its \emph{axiological and vitalistic
commitments}. The result is a threefold thesis:
\begin{enumerate}[label=(\roman*),nosep]
  \item \textbf{Informationalism:} an Idea is information, not a living
        substance.
  \item \textbf{Axiological Neutrality:} an Idea is neither good nor
        evil, neither true nor false in itself.
  \item \textbf{Domain-Absolute Independence:} an Idea does not depend
        on any mind, human or divine; it may, however, depend on the
        constitutive transformations of its domain.
\end{enumerate}

The three theses are logically independent: one may accept any subset
and reject the rest. Their conjunction, however, yields a distinctive
ontology whose consequences are worked out below. The document is a
sketch, not a closed system: where a transition is heuristic rather
than deductive, this is flagged.

Three terminological clarifications are in order at the outset, because
they will do substantial work later. First, the terms \emph{eternal},
\emph{immutable}, and \emph{aseitic} (self-existent, dependent on
nothing) are not synonyms: a structure may be immutable without being
eternal (a timeless abstractum that nonetheless depends on a divine
decree), and eternal without being aseitic (an everlasting emanation).
Second, the distinction between \emph{description} and
\emph{structure} is not the distinction between a linguistic item and a
non-linguistic one; it is the distinction between an arbitrary
partition of a state-space and a partition that is invariant under the
admissible transformations of that space. Third, the word
\emph{information} is used here in a structural sense, not in the
Shannon sense; the two are distinguished in
Definition~\ref{def:structure} and Remark~\ref{rem:shannon}. All three
distinctions are made precise in \S2.

\section{Fundamental Definitions}

\begin{definition}[Structure and Information]\label{def:structure}
Let $S$ be a space of possible states, and let $G$ be a groupoid of
\emph{admissible transformations} of $S$
(Definition~\ref{def:admissible}). A \emph{structure} on $S$ is a
partition $\Pi = \{B_1, B_2, \dots\}$ of $S$ into disjoint blocks
whose union is $S$.

A structure $\Pi$ is \emph{$G$-invariant} if for every $g \in G$ and
every block $B \in \Pi$, the image $g(B)$ is contained in a block of
$\Pi$, i.e.\ there is $B' \in \Pi$ with $g(B) \subseteq B'$. When $g$
is invertible on $S$ (as it is for group-valued symmetries), containment
strengthens to equality $g(B) = B'$. Two structures $\Pi, \Pi'$ are
\emph{$G$-equivalent} if they induce the same partition modulo the
action of $G$.

\textbf{Structural information} is a $G$-equivalence class of
$G$-invariant structures: the pattern of distinctions that survives all
admissible transformations of the state-space. A structure that is not
$G$-invariant is a \emph{description}, not structural information. The
\emph{informational content} of a structure is the pattern of
distinctions it imposes, modulo the identifications forced by $G$.

Three levels must be distinguished:
\begin{enumerate}[label=(\alph*),nosep]
  \item the \textbf{structure} itself --- the pattern of distinctions;
  \item its \textbf{carrier} --- a physical, neural, or symbolic medium;
  \item its \textbf{representation} --- a cognitive or linguistic
        encoding.
\end{enumerate}
A structure is \emph{not identical} to any one of its carriers or
representations. The information-theoretic Idea is a structure in sense
(a), and only in sense (a). The invariance condition is what
distinguishes structural information from a merely stipulated
partition: a partition that no admissible transformation preserves
carries no structural information about the system it partitions.
\end{definition}

\begin{remark}[Structural vs.\ Shannon information]\label{rem:shannon}
The term \emph{information} in this paper denotes
\emph{structural} information: a $G$-equivalence class of invariant
partitions, i.e.\ a pattern of distinctions robust under the admissible
transformations of a domain. It is \emph{not} Shannon information,
which is a measure of uncertainty reduction on a probability space and
is relative to a distribution over messages. The two notions are
compatible --- a Shannon channel can transmit a structural partition
--- but they answer different questions. Shannon information is
quantitative and probabilistic; structural information is qualitative
and invariance-theoretic. The choice of the structural notion is
deliberate: it is the notion that supports identity, embodiment,
representation, and misrepresentation, which are the notions the
ontology requires. For the Shannon notion and its scope, see
Cover and Thomas~\cite{CoverThomas2006}.
\end{remark}

\begin{definition}[Structural Identity, Embodiment, Representation]
\label{def:identity}
Two structures $\Pi$ and $\Pi'$ on $S$ are \emph{identical} if and only
if they induce the same partition: for every $s \in S$ and every block
$B \in \Pi$, there is a block $B' \in \Pi'$ with $s \in B
\Leftrightarrow s \in B'$. They are \emph{distinct} otherwise.

A structure $\Pi$ is \emph{embodied} in a carrier $C$ when there is a
map $f : C \to S$ such that, for every block $B \in \Pi$, the preimage
$f^{-1}(B)$ is a determinate subset of $C$. When $f$ is injective on
the relevant distinctions, the embodiment is \emph{faithful}.

A \emph{representation} of $\Pi$ is a pair $(R, g)$ where $g : S \to R$
is constant on each block of $\Pi$. Such a $g$ is a
\emph{homomorphism} from $\Pi$ onto the partition of $R$ induced by its
image. The representation is \emph{adequate} when $g$ is injective on
blocks, and \emph{coarsening} when distinct blocks may be identified.
Identity of structures is thus independent of the identity of their
carriers or representations.

Finally, a representation $(R, g)$ is \emph{misrepresenting} $\Pi$ when
there exist $s, s' \in S$ in distinct blocks of $\Pi$ with
$g(s) = g(s')$. This is the formal locus of error.
\end{definition}

\begin{definition}[Admissible Transformations]\label{def:admissible}
The groupoid $G$ of admissible transformations of $S$ is not chosen by
the observer. It is fixed by the domain: for a physical system, by its
symmetries (permutations of indistinguishable particles, changes of
inertial frame, gauge transformations); for a mathematical object, by
the equivalences internal to its category of presentation; for an
agent's practical domain, by the transformations that preserve the
conditions of agency (\S\ref{sec:good}). The choice of $G$
is a theoretical matter internal to the domain, not a psychological
matter. Objectivity is thus \emph{domain-relative without being
mind-relative}.

The phrase ``fixed by the domain'' is to be read as follows: for a
given domain, the admissible transformations are the morphisms of the
category in which that domain is presented. The choice of the category
is itself a theoretical choice, revisable in the light of the domain's
own demands; but it is not a choice made \emph{by} an observer
\emph{for} an observer. This is the sense in which Dummett's
distinction between realism and anti-realism about a domain turns on
what the domain itself requires, not on our epistemic convenience
\cite{Dummett1978}. See also Shapiro~\cite{Shapiro1997} for the
mathematical case.
\end{definition}

\begin{definition}[Information-Theoretic Idea]\label{def:idea}
An \textbf{Idea} (or \emph{information-theoretic eidos}) is a
structural information class in the sense of
Definition~\ref{def:structure} that is objective (invariant under the
admissible transformations of its domain), eternal, immutable, and
aseitic.
\end{definition}

\begin{definition}[Eternity, Immutability, Aseity]\label{def:modal}
A structure is \emph{eternal} if it is not located in time: there is no
time at which it begins or ceases to exist.

A structure is \emph{immutable} if no admissible transformation of its
state-space alters its partition. For a $G$-invariant structure this
condition is met by construction; immutability is therefore not an
independent modal predicate but a restatement of invariance. It is
retained here to mark what is \emph{excluded}: dependence on temporal
change.

A structure is \emph{aseitic} if neither its existence nor its identity
is caused, sustained, or constituted by any entity distinct from it.
Aseity is relative to a \emph{domain}: the structure's identity
conditions are exhausted by the domain-internal criteria of identity
(Definition~\ref{def:domain-abs}).
\end{definition}

\begin{definition}[Human Marker]\label{def:marker}
A \textbf{human marker} is a predicate --- such as \emph{true},
\emph{false}, \emph{good}, or \emph{evil} --- that a cognitive agent
applies to its \emph{relation} with an Idea. A marker is not an
intrinsic property of the Idea, and its application does not alter the
Idea.
\end{definition}

\begin{definition}[Marking and Pragmatic Commitment]\label{def:marking}
\textbf{Marking} is the act by which an agent applies a human marker to
its relation with an Idea. A marking is a \emph{pragmatic commitment}
when the agent's performance of the act of marking carries practical
entailments: to mark a representation as ``true'' is to undertake to
revise it should adequacy fail; to mark an action as ``good'' is to
undertake to promote the structures that condition one's agency. These
commitments are not derived from the Idea, nor from the marker alone.
They are \emph{posited} as constitutive of rational engagement with a
structured relation, in the tradition of Sellars~\cite{Sellars1956} and
Brandom~\cite{Brandom1994}; the position does not claim to deduce them
from the ontology of structures. This is a \emph{constructivist}
commitment, and it is a substantive metaethical position, not a
neutral one. It is logically independent of the three theses
(i)--(iii): one may accept them and reject the constructivism, or
accept the constructivism on other ontological grounds.
\end{definition}

\begin{definition}[Domain-Absolute Independence]\label{def:domain-abs}
An entity is \textbf{domain-absolutely independent} if neither its
existence nor its identity depends on any mind, subject, language,
culture, or divine act. Domain-absolute independence is the conjunction
of \emph{aseity} (Definition~\ref{def:modal}) with \emph{objectivity}
in the sense of Definition~\ref{def:admissible}.

Independence is \emph{absolute with respect to minds}, but
\emph{internal to a domain}: the admissible transformations that fix a
structure's identity are those of its domain, not those of an observer.
Whether the domain itself is further grounded --- in sets, categories,
or something else --- is left open. This means that aseity, in the
sense employed here, is relative to the domain: the position claims
independence from minds, not independence from every possible ground.
The stronger claim, that Ideas depend on \emph{nothing whatever}, is
not made and is not required by the three theses. The distinction
between the weaker and the stronger claim is the distinction between
\emph{domain-absolute} and \emph{absolutely absolute} independence;
only the former is asserted.
\end{definition}

\begin{remark}[Why the domain is not a hidden mind]
\label{rem:domain-not-mind}
A natural worry: if aseity is relative to a domain, and the domain is
fixed by ``theoretical choice'', is the domain not simply a disguised
form of mind-dependence? The reply turns on what fixes the domain.
A domain is fixed by its \emph{constitutive transformations} --- the
morphisms of the category in which it is presented. For the domain of
physical symmetries, these are the Poincaré group, gauge groups, and
permutation groups; for the domain of a mathematical structure, the
automorphism group of the structure; for the domain of agency, the
transformations that preserve the conditions of agency. In each case,
the transformations are \emph{discovered} in the domain, not
\emph{invented} by the theorist. The theorist chooses a
\emph{presentation} of the domain; the domain fixes which
presentations are adequate. This is the same distinction Dummett draws
between a domain's own standards of correctness and our epistemic
access to them \cite{Dummett1978}, and it is the distinction
Shapiro~\cite{Shapiro1997} draws between a structure and its
presentations. The domain is therefore not a hidden mind: it is the
locus of the constitutive transformations, and those transformations
are the domain's, not ours.
\end{remark}

\begin{remark}
The formal apparatus of Definitions~\ref{def:structure}--\ref{def:admissible}
is deliberately minimal. A partition is the weakest structure that
supports the notions of identity, embodiment, representation, and
misrepresentation; richer structures (orderings, algebras, categories)
can be recovered as partitions with additional constraints. The point
is not to reduce Ideas to set theory, but to make ``structural
determination'' precise enough to bear the arguments that follow.
\end{remark}

\section{The Three Theses}

\begin{claim}[Informationalism]
Ideas are information, not living beings.
\end{claim}

Consider the number $\pi$. It is eternal, immutable, and objective. It
does not grow, feel, desire, or act. It is not alive. Yet it is not
nothing: it is a determinate structure of distinctions --- a specific
position in the space of possible mathematical objects. The same
analysis applies to non-mathematical cases. The structure ``being a
cause'' is a determinate pattern of dependence among events; it does
not act, desire, or feel. The structure ``being a just procedure'' is
a determinate pattern of constraints on the distribution of burdens
and benefits --- constraints stated purely relationally, without appeal
to any evaluative primitive such as desert. It does not itself act
justly or unjustly. In each case, what we grasp is a structure, not a
living substance. By analogy with $\pi$, every Idea is of the same
ontological kind. Hence an Idea is information, not a living
substance.

\begin{remark}[The analogy, made precise]
The analogy with $\pi$ is not merely rhetorical. $\pi$ is a position
in the structure of the real numbers, and the real numbers form a
structure whose admissible transformations are the field automorphisms
fixing $\mathbb{Q}$. $\pi$ is invariant under these: any field
automorphism of $\mathbb{R}$ fixing $\mathbb{Q}$ fixes $\pi$, since
$\pi$ is definable from rationals by a first-order formula. The
analogy with Ideas is therefore an instance of the general
invariance criterion of Definition~\ref{def:structure}: an Idea is a
position in a structure that is invariant under the structure's
admissible transformations. What distinguishes an Idea from an
arbitrary mathematical object is not the kind of invariance but the
modal profile (eternal, immutable, aseitic).
\end{remark}

\begin{claim}[Axiological Neutrality]
Ideas are neutral with respect to truth, falsehood, good, and evil.
\end{claim}

Truth and falsehood are properties of propositions --- that is, of
human judgments \emph{about} an Idea. Good and evil are properties of
human attitudes and actions. An Idea, being neither a proposition nor
an action, cannot bear these predicates intrinsically.

``Neutrality'' here means the \emph{absence} of an axiological
predicate, not a predicate with a neutral value. An Idea is not
``neither good nor evil'' in the way a lukewarm action is; rather, the
question of good and evil does not arise for an Idea at all.

\begin{claim}[Domain-Absolute Independence]
Ideas are domain-absolutely independent of all minds.
\end{claim}

The argument proceeds in three steps, corresponding to the three
distinctions introduced in Definition~\ref{def:modal}.

\emph{First (immutability).} By Definition~\ref{def:identity}, the
identity of a structure is fixed by its partition of the space of
possible states, not by any mind's grasp of that partition. Two minds
may grasp the same structure, or fail to grasp it, without altering the
partition. Hence the identity of an Idea does not depend on any mind.
This establishes independence of \emph{identity} from minds.

\emph{Second (eternity).} An Idea is by definition not located in time.
A mind that created or sustained an Idea would have to stand in a
temporal relation to it --- a before and an after. Since an Idea admits
no such relation, no mind can create or sustain it. This establishes
independence of \emph{existence} from temporal modes of dependence.

\emph{Third (aseity).} It remains to exclude non-temporal modes of
dependence, such as the dependence of a thought on a thinker or of a
proposition on a language. Here the argument is structural: the
identity of a structure is fixed by its partition relative to the
admissible transformations of its domain
(Definition~\ref{def:admissible}). No appeal to a mind, a language, or
a divine act enters into this fixing. This is not a thesis about the
nature of sets; it is a thesis about the \emph{domain-internal}
criteria of identity. A structure is aseitic in the relevant sense
when its identity conditions are exhausted by these domain-internal
criteria. A mind may \emph{entertain} an Idea, but entertaining is a
relation to a structure whose identity is already fixed; the relation
is external to the relatum. This establishes independence of
\emph{existence} from non-temporal modes of dependence on minds.

The third step is the one most likely to be challenged, and it is
worth stating the challenge precisely. The challenge is that
``domain-internal criteria of identity'' are themselves
mind-dependent, because the domain is a theoretical construction. The
reply is the one given in Remark~\ref{rem:domain-not-mind}: the domain
is fixed by its constitutive transformations, and those
transformations are discovered in the domain, not invented by the
theorist. The theorist's contribution is the \emph{presentation}, not
the \emph{structure}. This is a substantive claim, and it is the
central claim of the paper; it is defended, not assumed, but it is
also not deduced from the definitions alone. The position is
therefore conditional on the correctness of this claim, and this is
flagged as an open problem in \S\ref{sec:limitations}.

From the independence of identity and the independence of existence
from minds, it follows that Ideas are domain-absolutely independent in
the sense of Definition~\ref{def:domain-abs}.

\begin{remark}
The three steps are not redundant. A critic who accepts the first but
rejects the second can hold that Ideas are mind-independent in content
but created in being; a critic who accepts the first two but rejects
the third can hold that Ideas are eternal, immutable, and yet sustained
by a non-temporal divine act. Information Idealism rejects both
positions, but it does so argument by argument, not by conflating the
three notions.
\end{remark}

\begin{remark}
A \emph{judgment} about an Idea may be false; the Idea itself is
neither true nor false. To say that a judgment is ``false for a
human'' is to predicate falsity of the judgment, not of the Idea. The
Idea's ontological status is unaffected by any human marker.
\end{remark}

\begin{remark}
Domain-absolute independence is an ontological claim, not an
epistemological one. How a finite mind can \emph{grasp} a
domain-absolutely independent Idea is a separate question, addressed
in \S\ref{sec:access}.
\end{remark}

\section{How Neutral Ideas Make Marking Possible}

A natural objection: if Ideas are wholly neutral and wholly
independent, why should a cognitive agent mark \emph{this} Idea as true
or good rather than any other? If the marker is wholly arbitrary, the
position loses explanatory power.

\subsection{Adequacy as a Structural Relation}\label{sec:access}

The marker is not arbitrary, because the relation between agent and
Idea is itself structured. When an agent's representation preserves
the partition of the Idea (Definition~\ref{def:identity}), the agent
is in a relation of \emph{adequacy} to it --- precisely the relation
of homomorphism described there. When the representation fails to
preserve the partition, the relation is one of \emph{inadequacy}. The
marker ``true'' is applied to the first relation, ``false'' to the
second. The marker is thus grounded in a structural fact about the
relation, not in a caprice of the agent.

The Idea contributes to this relation without itself bearing the
marker. Just as a key contributes to the act of opening a lock without
itself being ``open'' or ``closed'', an Idea contributes to the
adequacy or inadequacy of a cognitive relation without itself being
``true'' or ``false''. The axiological predicate attaches to the
relation, not to the relatum.

\subsection{A Theory of Misrepresentation}\label{sec:misrep}

The account so far explains how adequacy is possible, but not how
error is. Definition~\ref{def:identity} supplies the missing piece: a
representation misrepresents $\Pi$ when it identifies states that
$\Pi$ distinguishes. Error is thus not a mysterious failure of contact
with an Idea, nor a privation of being; it is a determinate structural
fact about a representation --- namely, failure of injectivity on
blocks. Three consequences follow.

\begin{enumerate}[label=(\alph*),nosep]
  \item \emph{Error is possible without the Idea being ``false''.} The
        Idea remains neutral; the falsity belongs to the relation.
  \item \emph{Error is detectable in principle.} Because the Idea's
        partition is fixed independently of any representation, a
        representation can be compared with it along the dimensions
        the admissible transformations make available. This is what
        ``checking'' means.
  \item \emph{Error is not total.} A coarsening representation
        preserves some distinctions while losing others; a
        misrepresentation preserves the partition on some region of
        $S$. Global misrepresentation is the degenerate case, not the
        norm.
\end{enumerate}

\subsection{Good and Evil as Markers of Agency}\label{sec:good}

The same analysis applies, \emph{mutatis mutandis}, to good and evil.
Let $A$ be an agent, and let $S_A$ be the state-space of the
conditions of $A$'s agency. The \emph{conditions of agency} are the
structural preconditions under which $A$ can form intentions,
deliberate, and execute actions: roughly, the capacity to represent
alternatives, to select among them on the basis of reasons, and to
translate selection into behaviour. The admissible transformations of
$S_A$ are those that preserve these conditions --- for instance,
transformations that do not destroy the agent's ability to form and
execute intentions. A structure $\Pi_A$ on $S_A$ is an Idea in the
relevant sense when it is invariant under these transformations.

An action is marked ``good'' when it promotes the structures that
condition $A$'s agency; it is marked ``evil'' when it obstructs them.
The Idea remains neutral; the marker remains a human one. This is a
\emph{constructivist} criterion: the good is not read off the Idea,
but from the conditions of agency as such. The choice of agency as the
anchoring notion is a normative starting point, not a deduction from
the ontology; other starting points are possible, and the position is
compatible with them. It is worth emphasizing that this is an
\emph{additional} thesis, logically independent of (i)--(iii): one may
accept the ontology of Ideas and reject the agency-anchored account of
good and evil, substituting another anchor (e.g.\ flourishing,
autonomy, or a contractarian procedure) without revising the ontology.

\begin{example}[A worked case: permutation symmetry]\label{sec:example}
To make the apparatus concrete, consider a domain consisting of three
indistinguishable particles, each of which may occupy one of two
states $\{0, 1\}$. The state-space is $S = \{0,1\}^3$, with $|S| = 8$.
The admissible transformations are the permutations of the three
particles: $G = S_3$, acting on $S$ by permuting coordinates.

The orbits of $G$ on $S$ are:
\begin{itemize}[nosep]
  \item $B_0 = \{(0,0,0)\}$ --- all particles in state $0$;
  \item $B_1 = \{(1,0,0), (0,1,0), (0,0,1)\}$ --- exactly one particle
        in state $1$;
  \item $B_2 = \{(1,1,0), (1,0,1), (0,1,1)\}$ --- exactly two
        particles in state $1$;
  \item $B_3 = \{(1,1,1)\}$ --- all particles in state $1$.
\end{itemize}
The partition $\Pi = \{B_0, B_1, B_2, B_3\}$ is the orbit partition;
it is $G$-invariant by construction, and it is the \emph{only}
$G$-invariant partition of $S$ that separates states by particle
count. It is, in the sense of Definition~\ref{def:structure}, a
structural information class --- the pattern of distinctions that
survives all permutations of the particles.

A representation of $\Pi$ is a map $g : S \to R$ constant on each
$B_i$. For instance, $g(s) = $ ``number of particles in state $1$''
is an adequate representation, with $R = \{0,1,2,3\}$. The map
$h(s) = $ ``parity of the number of particles in state $1$'' is a
coarsening representation, with $R = \{0,1\}$: it identifies $B_0$
with $B_2$ and $B_1$ with $B_3$. It is not misrepresenting, because
it does not identify states within a single block; it merely loses
distinctions. A map $k(s)$ that assigns the same value to $(0,0,0)$
and $(1,0,0)$ \emph{is} misrepresenting, because these lie in
different blocks of $\Pi$.

This example shows how the apparatus of
Definitions~\ref{def:structure}--\ref{def:identity} works in a
concrete case: invariance under a group action, orbit partition,
adequate vs.\ coarsening vs.\ misrepresenting representations. The
same template applies to richer domains (field automorphisms for
mathematical structures, gauge transformations for physical ones,
agency-preserving transformations for practical ones), with the
appropriate groupoid in place of $S_3$.
\end{example}

\section{Objections and Replies}

\begin{objection}[The Epistemic Access Problem]
If Ideas are domain-absolutely independent of all minds, how can any
mind \emph{identify} a particular Idea? An entity that is wholly
independent of every cognitive relation threatens to become a Kantian
thing-in-itself: something whose existence is asserted but whose
content is forever inaccessible. On such a view, the
information-theoretic Idea would be epistemically empty.
\end{objection}

\noindent
\emph{Reply.} The objection conflates \emph{ontological independence}
with \emph{epistemic inaccessibility}. Independence is a claim about
what fixes the Idea's identity and existence; it is not a claim about
what a mind can come to know. A structure can be fully determinate in
itself --- its partition fixed independently of any grasp --- and yet
be \emph{available} to a mind whose representations come to preserve
that partition. On this view, knowledge of an Idea is not a partial
participation in the Idea, nor a construction of it, but a relation of
structural adequacy between a representation and an independent
structure. Access is thus not mysterious: it is the ordinary fact that
a finite pattern can come to mirror an independent pattern. What is
rejected is only the further Platonic claim that the mind must be
\emph{akin} to the Idea, or that the Idea must somehow \emph{contain}
its own knowability.

The theory of misrepresentation in \S\ref{sec:misrep} completes the
reply. A critic who worries that ``structural adequacy'' is vacuous
because any representation can be gerrymandered to fit must reckon
with the fact that adequacy is checked along the admissible
transformations of the domain, not along arbitrary re-descriptions.
The transformations are not chosen by the knower; they are the
domain's own.

\begin{objection}[The Normativity Problem]
If Ideas are axiologically neutral, what grounds the \emph{normative
force} of the markers ``true'' and ``good''? A marker grounded only in
a structural fact about adequacy seems to be a description, not a
norm. Yet we treat ``true'' and ``good'' as \emph{claims} on us, not
as mere reports of correspondence.
\end{objection}

\noindent
\emph{Reply.} The normative force belongs to the \emph{act of
marking}, not to the marker considered in isolation. To mark a
representation as ``true'' is to undertake a pragmatic commitment
(Definition~\ref{def:marking}) to its adequacy: to revise it if
adequacy fails, to rely on it in inference, to treat it as a
constraint on further representation. To mark an action as ``good'' is
to undertake a pragmatic commitment to its promotion of the conditions
of agency. These commitments are not derived from the Idea, nor from
the marker alone. They are \emph{posited} as constitutive of what it
is to engage rationally with a structured relation, in the tradition
of Sellars~\cite{Sellars1956} and Brandom~\cite{Brandom1994}.

This is why Information Idealism can accept a Humean fact/value
distinction without collapsing into emotivism. The \emph{content} of
the marker is structural: it records a fact about the relation. The
\emph{force} of the marker is practical: it is the agent's undertaking.
The Idea supplies the structure; the agent supplies the commitment.
The commitment is not arbitrary, because an agent who refuses it
ceases to engage rationally with the relation --- ceases, in the
limit, to be an agent in the relevant sense. But the necessity here is
\emph{pragmatic}, not \emph{deductive}: it is the necessity of a
practice, not of a proposition. The position is therefore not
\emph{metaethically} neutral: it is a constructivism about normative
force. This is stated explicitly and not disguised as a consequence of
axiological neutrality at the level of Ideas. And, as noted in
\S\ref{sec:good}, this constructivism is logically independent of the
three ontological theses.

\begin{objection}[The Triviality Problem]
If every determinate structure is an Idea, then the position
multiplies Ideas beyond necessity: the structure of a particular heap
of sand, the structure of a passing cloud, the structure of a random
number sequence --- all become eternal, immutable, mind-independent
entities. This seems ontologically extravagant.
\end{objection}

\noindent
\emph{Reply.} The charge of extravagance presupposes that the
existence of a structure must be \emph{postulated} case by case. But
on the information-idealist view, structures are not postulated; they
are the objective partitions that any possible state-space admits.
That a heap of sand has a determinate structure is not an additional
ontological commitment; it is a consequence of the fact that the heap
occupies a determinate position in the space of possible
configurations. What is distinctive about \emph{Ideas}, in the strict
sense, is not that they are structures, but that they are structures
of a certain kind: those whose partitions are invariant under the
admissible transformations of their domain, and which are in addition
eternal, immutable, and aseitic. The ontology of structures is broad;
the ontology of \emph{Ideas} is narrow. The three theses apply to
Ideas so restricted, and the objection loses its force.

\begin{remark}[The restriction, made categorical]
\label{rem:categorical}
The restriction just invoked can be stated in a categorical idiom that
answers the demands of mathematically trained readers. Let
$\mathbf{S}$ be a \emph{groupoid} whose objects are the states of a
domain and whose morphisms are the admissible transformations
(Definition~\ref{def:admissible}). A \emph{structure} on
$\mathbf{S}$ is a functor $F : \mathbf{S} \to \mathbf{Set}$ that is
constant on isomorphism classes of objects --- i.e.\ $F(s) = F(s')$
whenever there is a morphism $s \to s'$ in $\mathbf{S}$. Such a
functor is precisely a partition of $\mathrm{Ob}(\mathbf{S})$
invariant under the morphisms of $\mathbf{S}$: the blocks are the
fibres $F^{-1}(x)$ for $x \in \mathrm{im}(F)$.

A \emph{representation} of a structure $F$ is a natural transformation
$\eta : F \Rightarrow G$ into another such functor $G$. The
representation is \emph{adequate} when $\eta$ is injective on each
fibre of $F$ (i.e.\ $\eta_s$ is injective for every object $s$), and
\emph{coarsening} when it is not. It is \emph{misrepresenting} when
there are objects $s, s'$ in distinct fibres of $F$ with
$\eta_s(F(s)) = \eta_{s'}(F(s'))$. This is the functorial form of
Definition~\ref{def:identity}.

A structure $F$ is an \emph{Idea} when, in addition to being a
structure in this sense, it is eternal, immutable, and aseitic. The
categorical formulation has three advantages over the group-theoretic
one. First, it accommodates partial symmetries (groupoids, not
groups). Second, it makes the notion of \emph{representation}
functorial, and adequacy a condition on natural transformations.
Third, it makes the \emph{domain-relativity} of objectivity precise
without making it mind-relative: different domains are different
groupoids, and the admissible transformations are the morphisms of the
relevant groupoid, not the choices of an observer. For the categorical
background, see Mac Lane~\cite{MacLane1998}; for the groupoid
formulation of symmetry in physics, see
Weinstein~\cite{Weinstein1996}.
\end{remark}

\section{Comparison with Classical Platonism}

The three theses can be read as a systematic revision of Plato's
positions. The table below summarizes the contrasts.

\begin{center}
\renewcommand{\arraystretch}{1.25}
\begin{tabular}{p{3.2cm} p{4.6cm} p{5.6cm}}
\toprule
\textbf{Aspect} & \textbf{Plato} & \textbf{Information Idealism} \\
\midrule
Nature of eidos & Living, intelligible substance & Information (structural determination) \\
Axiology & Good is the summit; evil is privation & No axiological predicate applies to Ideas; markers belong to relations \\
Truth / falsity & Ideas are truly being (\emph{ont\=os on}); falsity is a defect of cognition & Ideas are neutral; truth/falsity belong to judgments, grounded in structural adequacy; error is failure of injectivity on blocks \\
Independence & Ontologically independent, epistemically tied to the soul & Domain-absolutely independent: objective, eternal, immutable, aseitic \\
Access & Recollection; the soul is akin to the Ideas & Open: causal contact, structural isomorphism, or rational insight; adequacy is functorial \\
\bottomrule
\end{tabular}
\end{center}

\section{Comparison with Structural Realism}

Information Idealism is closest to \emph{structural realism} in the
philosophy of science and \emph{structuralism} in the philosophy of
mathematics. The comparison is worth making explicit, because it
locates the position's originality precisely.

Structural realists (Worrall~\cite{Worrall1989},
Ladyman~\cite{Ladyman1998}, French~\cite{French2014}) hold that what
is preserved across theory-change is structure, not intrinsic nature;
what exists are structures and the relations they articulate.
Mathematical structuralists (Shapiro~\cite{Shapiro1997},
Resnik~\cite{Resnik1997}, Hellman~\cite{Hellman1989}) hold that
mathematical objects are positions in structures, not independently
existing entities. Both positions are mind-independent in their
ontological commitments.

Information Idealism agrees with both on the primacy of structure and
on mind-independence, but differs on three points.

\begin{enumerate}[label=(\alph*),nosep]
  \item \emph{Scope.} Structural realism is a thesis about the
        \emph{physical} world; structuralism is a thesis about
        \emph{mathematical} objects. Information Idealism is a thesis
        about the ontology of eidos as such, and so ranges over
        physical, mathematical, and practical structures alike. It is
        a generalization of both, not a competitor.
  \item \emph{Axiological neutrality.} Structural realists and
        structuralists are typically silent on the relation between
        structure and value. Information Idealism makes axiological
        neutrality a central thesis, and supplies an account of
        markers (\S4) that neither tradition offers.
  \item \emph{The modal triad.} Structural realists and
        structuralists typically assert mind-independence without
        distinguishing objectivity, eternity, immutability, and
        aseity. Information Idealism separates them
        (Definition~\ref{def:modal}) and argues for each separately
        (Claim 3.3). This is what makes the position
        \emph{domain-absolutely} independent in a sense stronger than
        either tradition usually claims, while stopping short of the
        stronger claim of independence from every possible ground.
\end{enumerate}

The originality of Information Idealism thus lies not in the primacy
of structure --- that is shared with structural realism and
structuralism --- but in the \emph{combination} of structural ontology
with explicit axiological neutrality and a differentiated account of
domain-absolute independence. Each ingredient is traditional; the
conjunction is not.

\section{Limitations and Open Problems}\label{sec:limitations}

The position is programmatic, and it is worth stating the open
problems explicitly, both to forestall the charge of overclaiming and
to indicate directions for further work.

\begin{enumerate}[label=(\arabic*),nosep]
  \item \emph{The grounding of the domain.} The central open problem
        is the one identified in Remark~\ref{rem:domain-not-mind}:
        what fixes a domain's constitutive transformations, and why
        is this fixing not itself a form of mind-dependence? The
        reply offered --- that the transformations are discovered,
        not invented --- is substantive but not fully argued. A
        fully worked-out account would need to say, for each kind of
        domain (physical, mathematical, practical), what it is for a
        transformation to be constitutive, and how we can be wrong
        about it. This is the most important direction for further
        work.
  \item \emph{The identification of Ideas across domains.} If two
        domains have different constitutive transformations, can they
        share an Idea? The position as stated leaves this open. A
        natural answer is that cross-domain identity is
        \emph{functorial}: an Idea is shared when there is a functor
        between the groupoids that preserves the relevant partitions.
        But this answer needs to be made precise and defended.
  \item \emph{The completeness of the modal triad.} The triad
        (eternal, immutable, aseitic) is offered as a
        \emph{differentiated} account of independence. It is not
        claimed to be exhaustive. Other modal notions --- necessity,
        essentiality, fundamentality --- might also be relevant, and
        their relation to the triad is not spelled out.
  \item \emph{The status of the constructivism.} The constructivist
        account of normative force is logically independent of the
        ontology, but it is not for that reason optional: if it
        fails, the position loses its account of the normativity of
        ``true'' and ``good''. The account is stated in the tradition
        of Sellars and Brandom, but the details --- in particular, the
        account of the pragmatic commitments involved in marking ---
        are left schematic.
  \item \emph{The epistemology of adequacy.} The account of
        misrepresentation as failure of injectivity is formal and
        clear, but the account of how agents \emph{come to know}
        whether their representations are adequate is not. The
        position says that adequacy is checked along the admissible
        transformations of the domain, but it does not say how the
        agent identifies the relevant transformations in practice.
        This is a problem for the epistemology, not the ontology, but
        it is a problem.
\end{enumerate}

\section{Conclusion}

Information Idealism offers a coherent and modern reformulation of
Platonism. It preserves the realist core of Plato's theory while
discarding its vitalistic and moralistic overtones. An Idea is
information --- objective structural determination; it is neutral in
the axiological sense; it is domain-absolutely independent, in the
differentiated sense of objective, eternal, immutable, and aseitic.
Truth, falsehood, good, and evil are not properties of Ideas but
markers of the human relation to them. The formal apparatus ---
partitions, admissible transformations, functorial representations,
misrepresentation as failure of injectivity --- is minimal but
sufficient to state the position precisely, to answer the standard
objections, and to locate the view in relation to classical Platonism
and to structural realism.

Information Idealism is offered as a programmatic sketch: the formal
core is minimal, several transitions are argued heuristically, and the
normative theory is constructivist rather than derived. Its aim is to
state the position clearly enough to be criticized. The open problems
of \S\ref{sec:limitations} indicate the directions in which the
position would need to be developed to be more than a sketch.

\begin{thebibliography}{99}

\bibitem{Brandom1994}
R.~Brandom.
\newblock \emph{Making It Explicit: Reasoning, Representing, and
Discursive Commitment}.
\newblock Harvard University Press, 1994.

\bibitem{CoverThomas2006}
T.~M. Cover and J.~A. Thomas.
\newblock \emph{Elements of Information Theory}.
\newblock Wiley-Interscience, 2nd edition, 2006.

\bibitem{Dummett1978}
M.~Dummett.
\newblock \emph{Truth and Other Enigmas}.
\newblock Harvard University Press, 1978.

\bibitem{French2014}
S.~French.
\newblock \emph{The Structure of the World: Metaphysics and
Representation}.
\newblock Oxford University Press, 2014.

\bibitem{Hellman1989}
G.~Hellman.
\newblock \emph{Mathematics without Numbers: Towards a Modal-Structural
Interpretation}.
\newblock Clarendon Press, 1989.

\bibitem{Ladyman1998}
J.~Ladyman.
\newblock What is structural realism?
\newblock \emph{Studies in History and Philosophy of Science},
29(3):409--424, 1998.

\bibitem{MacLane1998}
S.~Mac Lane.
\newblock \emph{Categories for the Working Mathematician}.
\newblock Springer, 2nd edition, 1998.

\bibitem{Resnik1997}
M.~D. Resnik.
\newblock \emph{Mathematics as a Science of Patterns}.
\newblock Clarendon Press, 1997.

\bibitem{Sellars1956}
W.~Sellars.
\newblock Empiricism and the philosophy of mind.
\newblock In H.~Feigl and M.~Scriven, editors, \emph{Minnesota Studies
in the Philosophy of Science}, volume~1, pages 253--329. University of
Minnesota Press, 1956.

\bibitem{Shapiro1997}
S.~Shapiro.
\newblock \emph{Philosophy of Mathematics: Structure and Ontology}.
\newblock Oxford University Press, 1997.

\bibitem{Weinstein1996}
A.~Weinstein.
\newblock Groupoids: unifying internal and external symmetry.
\newblock \emph{Notices of the American Mathematical Society},
43(7):744--752, 1996.

\bibitem{Worrall1989}
J.~Worrall.
\newblock Structural realism: the best of both worlds?
\newblock \emph{Dialectica}, 43(1--2):99--124, 1989.

\end{thebibliography}

\end{document}